% Generated by task tex-docs from dossiers/docs-debt.md, section C (vocabulary). Do not edit by hand.
\section{Vocabulary of the documents}\label{sec:docs-vocabulary}

Counts are whole-word, case-insensitive occurrences (\texttt{vocab\_\allowbreak{}counts.\allowbreak{}py}; BP, ST, TB, HC as in \cref{sec:docs-inventory}). ``Keep'' follows the brief's vocabulary where it decides.

{\footnotesize\setlength{\tabcolsep}{3pt}
\begin{longtable}{@{}>{\raggedright\arraybackslash}p{0.143\linewidth}>{\raggedright\arraybackslash}p{0.316\linewidth}>{\raggedright\arraybackslash}p{0.210\linewidth}>{\raggedright\arraybackslash}p{0.287\linewidth}@{}}
\caption{Variants of one term, where each is used, and the term to keep.}\label{tab:docs-vocabulary}\\
\toprule
\textbf{One thing} & \textbf{Variants, with counts} & \textbf{Keep} & \textbf{Remark} \\
\midrule
\endfirsthead
\toprule
\textbf{One thing} & \textbf{Variants, with counts} & \textbf{Keep} & \textbf{Remark} \\
\midrule
\endhead
\bottomrule
\endfoot
The document that says what is wanted & ``roadmap'': BP 25, ST 16, TB 13; ``blueprint'': BP 5, ST 3, TB 3, HC 11, the handoffs 37. The file is \texttt{protocol-\allowbreak{}blueprint.\allowbreak{}md}; its title is ``Roadmap: ...'' & blueprint & a reader of the tracker meets both words for one file in one sentence (\texttt{issue-\allowbreak{}12-\allowbreak{}body.\allowbreak{}md:\allowbreak{}1}) \\\addlinespace[2pt]
The tracking issue & ``dashboard'': BP 7, ST 2; ``tracker'': BP 2; ``tracking issue'': BP 1, ST 2; ``this issue'' & tracker & at \texttt{protocol-\allowbreak{}blueprint.\allowbreak{}md:\allowbreak{}1489} ``the tracker'' means all of GitHub's issues, not issue 12 \\\addlinespace[2pt]
A unit of work & ``hole'': BP 48, TB 29, HC 21; ``layer'': BP 172, HC 40; ``unit'': BP 29; ``slice'': BP 9; ``half'': BP 10; ``leaf'': see below & ``hole'', defined once (``a unit of work the spine leaves open: a phase, a generic component, the adaptor, or a piece of the compilation''), always with its name; ``layer'' for the blueprint's section number & ``hole'' is not a word a zkVM engineer knows, so its definition goes in the index; a hole has up to four identifiers today (hole code, layer number, the closed issue's number, the pull request). The earlier draft proposed ``unit of work'' instead; the second pass keeps ``hole'', the brief's word, the word of the spine's Lean docstrings (``the hole interfaces'') and of 48 places in the blueprint. In \cref{sec:docs-proposal-owners} and \cref{sec:docs-proposal-codes} below, ``unit of work'' and ``unit'' mean a hole \\\addlinespace[2pt]
The table of upstream gaps & ``ledger'': BP 24, ST 10; ``upstream ledger'': BP 2; ``upstream watch'': BP 3, HC 5; ``upstream candidates'' (\texttt{leanth-\allowbreak{}reuse.\allowbreak{}md:\allowbreak{}447}) & ``upstream ledger'' for what is admitted or absent at the pin and what replaces it; ``upstream watch'' for external work in flight & two different tables; the blueprint's hole table has a third, the column ``Existing work'' \\\addlinespace[2pt]
The relation between two phases & ``seam'': BP 62; ``interface'': BP 28; ``boundary'': BP 17 & seam & ``interface'' also means an assumed structure (\texttt{FlockInterface}) and the list of public names; ``boundary'' also means the boundary blocks of the bus and the section ``Boundaries'' with other roadmaps \\\addlinespace[2pt]
One step of the protocol & ``phase'': BP 131; ``reduction'': BP 17; ``component'': BP 60; ``stage'': BP 3 & phase for the six steps; component for a generic piece & \texttt{Component} is three things: Clean's table circuit (\texttt{:\allowbreak{}291}), the spine's \texttt{Component.\allowbreak{}Def} (\texttt{:\allowbreak{}501}), ArkLib's \texttt{Proof\allowbreak{}System/\allowbreak{}Component/\allowbreak{}} (\texttt{:\allowbreak{}244}) \\\addlinespace[2pt]
The committed object & ``stack'': BP 41; ``stacked columns''; ``the oracle'': BP 51; ``committed column/polynomial''; ``\texttt{q}'' & the stack & ``column'' is both a table's column and the stack itself (\texttt{Column μ}) \\\addlinespace[2pt]
A coordinate of \texttt{E} over \texttt{K} & ``limb'': BP 24; ``lane'': BP 5 & limb & ``lane'' has three meanings: a word of the BLAKE2s block (\texttt{:\allowbreak{}325}), the four \texttt{K}-words of the public input in leanISA (\texttt{leanisa-\allowbreak{}blueprint.\allowbreak{}md:\allowbreak{}356}), and a track of work in the Flock roadmap (\texttt{:\allowbreak{}1446}). ``limb'' is also used for a column: ``the three memory limbs'' \\\addlinespace[2pt]
The claims handed to the opening & ``claim pool'': BP 5; ``pool'', ``pooled'': BP 19, RP 33 & claim pool & ``claim'' is also the act of taking a unit of work (``claimed'': BP 7, TB 7) \\\addlinespace[2pt]
Done & ``landed'': TB 13, ST 9; ``built'': BP 8, ST 5; ``on \texttt{main}'': ST 10; ``merged'': ST 16; ``met'' (of a finding): the handoffs 39; ``done''; ``taken'' (of a decision) & ``on main'' for code; ``met'' for a finding; ``taken'' for a decision & the tracker defines four states (\texttt{issue-\allowbreak{}12-\allowbreak{}body.\allowbreak{}md:\allowbreak{}18}); the status uses a fifth, ``built, draft'' \\\addlinespace[2pt]
The import rule & ``the wall'': BP 7, ST 4, RL 6 & ``the import rule of the proof system'', stated once & a metaphor; the rule is one sentence (\texttt{:\allowbreak{}331}) \\\addlinespace[2pt]
The map between the two relations & ``adaptor'': BP 29; ``adapter'': ST 1 (another thing, a compatibility file); ``bridge'': BP 11; ``refinement'': BP 13 & adaptor for the leanISA one; refinement for the generic notion & ``the Clean bridge'' is Layer 2, a different unit \\\addlinespace[2pt]
The protocol & ``oracle protocol'': BP 11; ``PIOP'' in names (\texttt{leanVmPiop}); ``IOR'': BP 1; ``IOPP'' (\texttt{leanVmIopp}): BP 5 & oracle protocol & \texttt{leanVmIopp} names as a proximity proof the compiled protocol ``WHIR in place of the evaluation oracle'' (\texttt{:\allowbreak{}1214-\allowbreak{}1215}) \\\addlinespace[2pt]
A bus entry & ``flush'': BP 12; ``interaction'': ST 10; ``tuple'' & flush, as the specification & ``side'' (\texttt{Side}, spine) against ``direction'' (\texttt{Direction}, Layer 2 sketch, and ArkLib's message direction) \\\addlinespace[2pt]
The public data & ``public input'': ST 23; ``public statement'': BP 3; ``public words'': BP 6; ``public line'': BP 4 & public input for the two words; public line for the spine's object & \texttt{PublicInput} is a leanISA type and also a namespace of the proof system (\texttt{Public\allowbreak{}Input.\allowbreak{}check}) \\\addlinespace[2pt]
The work item inside a layer & ``leaf'', ``leaves'': BP 26 & - & four meanings: a leaf of the product tree (\texttt{:\allowbreak{}319}), ``the leaves of Layer 1 already in flight'' (\texttt{:\allowbreak{}591}), ``the missing leaf'' of ArkLib's sumcheck (\texttt{:\allowbreak{}259}), ``packing leaves'' (\texttt{:\allowbreak{}267}) \\
\end{longtable}}

Terms a zkVM engineer would not recognize without the blueprint's definitions: hole, seam, spine, the wall, ledger, frontier, toy (instance), Category A and Category B, M3 (expanded only in the references, \texttt{:\allowbreak{}1539}), IOR. Of these only ``M3'' and ``IOR'' are terms of art elsewhere.

Words with two meanings that matter: \textbf{claim} (an evaluation claim; the act of taking work), \textbf{instance} (\texttt{M3Instance}; a Lean type-class instance; ``an instance of a theorem''), \textbf{statement} (the public statement \texttt{I.\allowbreak{}Stmt}; a theorem's statement; Clean's \texttt{Ensemble.\allowbreak{}Statement}), \textbf{Layer N} (the proof system's and leanISA's, both from 0; the table at \texttt{protocol-\allowbreak{}blueprint.\allowbreak{}md:\allowbreak{}215-\allowbreak{}222} uses the bare form for leanISA's and ``Layer 11 here'' for its own in the same row), \textbf{component}, \textbf{lane}, \textbf{leaf}, \textbf{Direction}.
